\chapter{总结与展望}

\section{全文总结}

本文针对复杂场景下物理仿真中视觉表征与物理属性分离的核心问题，设计并实现了物理知识嵌入的复杂场景生成与仿真系统。系统以文本或参考图像为输入，贯通了从三维场景构建、材质理解与映射、到多材质物理仿真与参数优化的完整链路。

在场景构建与材质初始化方面，系统通过 Director3D\cite{director3d} 扩散轨迹生成与 FlashWorld\cite{flashworld} 前馈重建构建三维高斯场景雏形，并引入 Depth-Anything-V2\cite{depthanythingv2} 单目深度估计与基于语义块的 K-Means\cite{kmeans} 深度聚类分层，在空间维度上解耦前景与背景。SAM\cite{sam} 自动分割提供语义块轮廓后，Qwen3-VL-32B\cite{qwen3vl} 视觉语言大模型对每个孤立区域进行逐层材质分类，通过 \texttt{material\_config.json} 经验库将语义标签映射为杨氏模量、泊松比与密度的初始物理参数范围，为后续物理参数优化提供约束先验。

在特征蒸馏方面，系统基于 Feature Splatting\cite{featuresplatting} 框架，将 CLIP\cite{clip} 与 DINOv2\cite{dinov2} 二维视觉特征通过可微 $\alpha$-blending 联合蒸馏至三维高斯原语，并通过独立的语义标签解码器将 VLM 材质分类结果同步映射至高斯空间。在此之上，引入五阶段前景分割管线——集成了得分引导 DBSCAN\cite{dbscan} 聚类、证据融合（CLIP 对象级/零件级相似度、SAM 掩码反投影、DINOv2 KNN 平滑）与连通分量过滤——实现前景主体与静态背景的解耦，避免大规模背景在 MPM 仿真中引发网格刚体化。

在物理仿真与参数优化方面，系统以物质点法（MPM）\cite{mpm,physgaussian} 为力学求解器，为不同材质区域的质点分配独立的本构模型（弹性、粘弹性、塑性等），并集成 DreamPhysics\cite{dreamphysics} 框架以 SVD\cite{svd} 作为隐式物理先验。通过分数蒸馏采样（SDS）将 MPM 仿真帧与 SVD 预期运动在潜空间中的分布偏差转化为弹性模量的优化梯度，经 Warp\cite{warp} 全可微物理管线反传至每一物质点。优化过程中，系统集成了运行时 CFL 自适应时间步长、逐粒子 CFL 钳位、体素材质正则化与按材质类型梯度归一化等多项数值稳定措施。

本文的核心贡献有三点：一是，构建了从文本/图像到可仿真虚拟场景的端到端生成式管线，连通了场景生成、材质理解与物理仿真三个环节；二是，以 VLM 驱动的自动材质识别与经验映射替代人工物理参数指定，为不同材质区域提供稳定的参数先验；三是，扩展前景—背景解耦仿真架构与视频先验引导的多材质物理参数梯度优化框架，在缺乏真值物理标注的条件下实现参数的自监督收敛纠偏。

\section{系统局限性分析}

尽管本文系统在多个场景案例上完成了从生成到仿真的完整链路，但是当前系统仍存在以下局限：

\textbf{生成场景的几何质量。}系统以 FlashWorld 作为前馈式场景生成骨干，其依赖 Director3D 生成的相机轨迹与观测序列。当输入文本描述过于复杂或包含长尾概念时，生成的相机轨迹可能遗漏关键观测角度，导致重建场景出现几何空洞、纹理拉伸或物体结构不完整等问题。此外，FlashWorld 的前馈特性使其无法像 3DGS\cite{3dgs} 原位训练那样对局部几何进行精细修正，在面对高曲率边界或薄壁结构时容易出现高斯退化。

\textbf{材质识别的粒度与精度。}当前 VLM 材质分类限定于 9 类预设材质库，在面对跨材质复合材料（如碳纤维增强聚合物、金属镀层玻璃）时缺乏细粒度的子类别区分能力。VLM 在区分视觉表观相近的材质时（如透明玻璃与透明塑料、抛光金属与镜面塑料）存在分类混淆。此外，单帧推理的不确定性使多数投票机制（3 次采样）仅能提供有限的置信度提升，尚未从根本上解决 VLM 的材质幻觉问题。

\textbf{VLM 推断到物理参数映射的精度。}\texttt{material\_config.json} 经验库提供的物理参数范围为宏观经验区间（如石材 $E\approx10^8$\,Pa），未考虑材质内部的微观结构与各向异性特征。对于硬度等级的判断（soft/medium/hard），VLM 仅依赖视觉表观信息，缺乏力学反馈信号的闭环验证，可能导致参数先验过于宽泛，为后续 SDS 优化增加搜索负担。

\textbf{仿真数值稳定性边界。}尽管系统引入 CFL 自适应时间步长与逐粒子钳位等措施，MPM 多材质耦合仿真在以下场景仍可能数值发散：(1) 杨氏模量跨材质差异过大时，高刚度材质迫使全局时间步长急剧缩小；(2) 近不可压缩材料（$\nu>0.4$）的波速放大效应进一步加剧 CFL 约束；(3) 粒子填充密度不足时，稀疏网格导致应力低分辨率传递，产生虚假的局部变形。此外，高斯壳体到 MPM 粒子的内部刚性填充策略仅适用于实体物体，对中空薄壳结构需额外几何推理。

\textbf{SDS 优化的收敛效率与物理合理性。}SDS 梯度具有高方差特性\cite{dreamfusion}，单次优化需数百次迭代才能在视觉上产生可分辨的形变差异，且收敛方向对时间步采样策略敏感。SVD\cite{svd} 视频先验虽蕴含运动常识，但其训练数据中的非物理运动（如慢动作镜头、人工特效）可能引入偏置，导致优化后的物理参数偏离力学真值。当前优化仅涉及杨氏模量 $E$ 的梯度更新，泊松比 $\nu$ 与密度 $\rho$ 保持经验赋值，限制了多参数联合优化的上界。

\textbf{前景分割的文本依赖性。}五阶段前景分割管线对文本查询的语义覆盖范围有较强依赖——当目标物体缺乏直观的文本描述词（如抽象雕塑、异形器具），CLIP 文本匹配难以提供有效得分，需退化为 SAM 掩码反投影的独立链路，后者缺乏语义筛选能力，可能将空间邻接的背景碎片误标为前景。

\section{未来研究展望}

基于上述局限性分析，未来研究工作可从以下几个方向展开：

\textbf{高质量场景生成。}探索更先进的前馈式场景生成模型或在重建后追加短时原位优化微调机制，在保持生成效率的同时提升几何精度。引入多模态提示条件（文本结合草图或参考样式）以增强对复杂语义场景的生成控制能力。

\textbf{细粒度物理材质感知。}扩展 VLM 材质分类体系至更丰富的材质库，引入多模态特征融合联合推断，并利用物理仿真结果的视觉反馈对 VLM 分类进行闭环校验与纠偏。探索基于学习的物理参数回归网络——以 VLM 特征与几何特征为输入、以 MPM 仿真误差为监督信号——替代经验范围映射，实现从视觉表观到力学参数的端到端估计。

\textbf{多材质仿真的鲁棒性增强。}在 MPM 求解器中引入自适应本构切换机制：当某区域参数更新后匹配另一本构模型的更优拟合条件时，自动迁移至新的本构关系，减少经验预设带来的模型偏差。探索粒子-网格自适应细化策略，在高形变梯度区域动态提升粒子分辨率以捕获细粒度力学响应。对中空结构进行拓扑推断与自适应粒子填充，使高斯壳体能够支撑准确的力学计算。

\textbf{改进的视频先验引导优化。}研究低方差梯度估计器以替代标准 SDS，加速收敛并降低优化不稳定性。拓展优化空间至多物理参数联合寻优（$E$、$\nu$、$\rho$ 及局部本构参数），利用更先进的视频生成模型提供的时空运动先验。引入具有真实物理偏置的训练集过滤策略，减少视频扩散模型中非物理运动的污染。

\textbf{多物体交互与复杂碰撞。}当前系统主要聚焦单一前景物体的受力形变，未来可扩展至多物体碰撞交互、刚柔耦合运动以及流体-固体双向耦合仿真。通过引入物体间相对位姿预测与接触面动力学约束，使系统能够模拟抓取、堆叠、倾倒等更丰富的物理交互场景，为机器人视觉-物理联合仿真铺平道路。

\textbf{从仿真到现实的物理属性迁移。}利用本文系统为真实物体快速构建可仿真数字孪生，通过少量真实物理交互数据对仿真参数进行域适应微调，缩小 Sim2Real 的物理属性鸿沟。该方向可将本系统的输出接入机器人仿真平台，为具身智能体的物理推理与操作规划提供高保真虚拟训练环境。
